\section{从内部评分到科学证据：AMR 的受控复现}

架构图说明系统如何产生假设，却没有回答“这些假设到底值不值得相信”。本章先建立证据阶梯，再分析抗菌药耐药性案例。这个案例的特殊价值在于：人类团队已经获得但尚未公开的答案，可以用来检测系统是否在没有目标答案可检索时独立重建同一机制。

\subsection{不要把七种证据压成一个 success rate}

科研 Agent 的证据至少分为内部 judge、文献一致性、专家评价、隐藏结果复现、前瞻实验、独立重复和临床证据。它们回答不同问题。一个机制可以文献合理却实验失败，也可以在细胞中有效却在人体中无效。课程后面的案例必须逐个标明处于哪一层。

\begin{importantbox}{科学证据阶梯}
\begin{tabularx}{\textwidth}{L{0.22\textwidth}X X}
\toprule
层级 & 能支持的结论 & 仍不能支持的结论 \\
\midrule
内部 judge / Elo & 候选在 rubric 下相对更好 & 真实、新颖、可复现 \\
文献与工具检查 & 与公开知识较一致、结构较可行 & 因果机制成立 \\
领域专家评价 & 候选值得进一步检查 & 实验一定成功 \\
隐藏结果复现 & 系统能重建未知答案 & 系统前瞻发现了新事实 \\
前瞻实验 & 某个预注册或事先提出的节点得到支持 & 普遍机制与临床获益 \\
独立重复 / 临床研究 & 结果跨团队或人群更稳健 & 无条件外推到所有环境 \\
\bottomrule
\end{tabularx}
\end{importantbox}

若实验容量有限，可以按期望信息增益分配资源：
$$
P(h_i)=\frac{\mathbb{E}[\Delta I_i]\,p_i}
{C_i+\lambda R_i}.
$$
其中：
\begin{itemize}
  \item $\mathbb{E}[\Delta I_i]$ 是实验对区分机制的期望信息增益；
  \item $p_i$ 是候选通过当前证据校准后的可信度；
  \item $C_i$ 是金钱、时间和样本成本；
  \item $R_i$ 是安全、伦理或不可逆风险，$\lambda$ 是风险权重。
\end{itemize}

\teachervoice{00:37:18--00:39:05，老师预判未来瓶颈可能从“没有好假设”转为“好像有太多好假设，验证不过来”。因此系统必须优化实验选择，而不是只最大化报告长度。}

\subsection{cf-PICI 问题：怎样跨物种传播}

Capsid-forming phage-inducible chromosomal islands（cf-PICIs）是依赖噬菌体传播的移动遗传元件。它们能形成自己的 capsid（装 DNA 的“头部”），却需要利用 helper phage 的 tail 完成感染与宿主识别。人类团队经过多年实验发现，cf-PICI 可借用多种噬菌体尾部，形成不同宿主范围的嵌合颗粒，从而解释相同 cf-PICI 为什么出现在多种细菌中。

\lecturefigure{slide-008.jpg}{AI co-scientist 对 AMR 相关基因转移机制的受控复现流程}{Stanford CS25 V6 Lecture 07 官方录像 00:40:20；bioRxiv:2025.02.19.639094v1}

\begin{knowledgebox}{读图：左边是多年实验，右边是数日假设搜索}
左侧从提出 cf-PICI、观察跨物种分布、构造多种机制，到实验确认新型传播路径，跨越约十年；右侧在不知道未公开答案的前提下，由 AI co-scientist 生成、排序和审查候选，最高排名假设复现“借用多种 phage tail 扩大 host range”的核心。中间相同研究问题是两条路线可比较的锚点。
\end{knowledgebox}

\begin{knowledgebox}{First-use glossary：cf-PICI 机制}
\begin{tabularx}{\textwidth}{L{0.22\textwidth}X X}
\toprule
术语 & 含义 & 在机制中的作用 \\
\midrule
Capsid & 包装遗传物质的蛋白质头部 & cf-PICI 可形成自己的 capsid \\
Phage tail & 负责识别宿主并注入遗传物质的尾部结构 & 被 cf-PICI 借用，决定可感染宿主 \\
Host range & 一个颗粒能进入哪些细菌宿主 & 多种 tail 扩展跨物种传播范围 \\
Recapitulation & 在不知道隐藏答案时重建已被另一团队发现的机制 & 检验假设生成能力，不是前瞻发现 \\
\bottomrule
\end{tabularx}
\end{knowledgebox}

\teachervoice{00:40:13--00:44:13，老师明确把这个结果称为 recapitulation：合作团队先在现实中完成发现，只是答案尚未公开。系统用数日重建了多年实验得到的核心解释。}

\begin{warningbox}{为什么这是强控制，却仍不是“AI 发现”}
强点在于目标答案不在公开文献中，降低了直接检索复述的可能；弱点在于研究问题、公开背景和候选评价仍由人类提供，且实验是合作团队事先完成的。最准确表述是“隐藏结果复现”或“平行 in-silico hypothesis generation”。
\end{warningbox}

\subsection{泄漏控制与时间版本}

判断隐藏结果复现是否可信，要检查模型知识截止、预印本时间、实验团队何时提交论文、提示中提供了哪些资料，以及评测者是否在看到答案后调整 rubric。本讲使用课堂日前的 bioRxiv v1 与 arXiv v1；2026-05-14 之后的 Nature 版本和后续修订不能被当成课堂已经拥有的证据。

\begin{warningbox}{Temporal leakage 检查表}
记录模型版本与运行日期；冻结输入文档；核查预印本、会议摘要、代码和社交媒体；让不知道目标答案的评审先评价候选；保留完整候选排名而不只展示命中项。若缺少这些记录，“未知答案”可能只是研究者不知道答案已进入训练或检索渠道。
\end{warningbox}

\subsection{本章小结}

AMR 案例说明 AI co-scientist 不只会把文献摘要改写成假设：在答案未公开的受控问题上，最高排名候选重建了多种 phage tail 扩展 cf-PICI host range 的核心机制。但这仍是隐藏结果复现，而不是前瞻发现。其真正意义是建立一层比内部 judge 更强、比独立前瞻实验更弱的证据，并暴露时间泄漏与版本冻结的重要性。

